\part{Cyclotomic limits and the Fabius--Rvachev deformation}

\section{The boundary and probability regimes are different}
\label{sec:cyclotomic-fabius-interface}

Two uses of a geometric parameter occur in this volume.  In
\cref{sec:ip-cyclotomic}, the base \(q\) approaches the unit circle and the
Lambert series develops resonant residue classes.  In the present part,
\(0<q<1\) is the contraction ratio of a random series.  The Rvachev member is
the perfectly regular interior value \(q=1/2\).  Thus the radial inverse
\cref{thm:ip-radial-cyclotomic} and the finite cyclotomic analysis of
\cref{sec:gaussian-cyclotomic} remain available when a product formula reaches
the boundary, but they should not be conflated with statistical recovery of
the geometric contraction parameter.

The law, support, CDF refinement, smooth density, and half-base
identification are represented in Lean by the modules
\lean{FabiusFunction.GeometricUniformLaw},
\lean{FabiusFunction.GeometricUniformCDF}, and
\lean{FabiusFunction.ProbabilityRepresentation}.  In particular,
\lean{geometricUniformCDF_one_half_eq_fabiusReal} is the exact formal bridge
at \(q=1/2\).  The phase-bearing sinc-product bridge is formalized separately
in \lean{FabiusFunction.GeometricSincCharacteristicFunction}, as
\lean{charFun_geometricUniformDistribution_eq_phase_mul_geometricSincProduct}.
The arbitrary-power Fourier decay in \cref{eq:fr-fourier-rapid}, its
parameter-derivative strengthening in \cref{eq:fr-q-derivative-rapid}, and the
cumulant inversions later in this chapter are complete human proofs at the
frontier level; they are not claimed here as existing Lean declarations.

\section{The geometric uniform law and its sinc product}
\label{sec:fr-geometric-law}

Let \(V_0,V_1,\ldots\) be independent random variables uniformly distributed
on \([0,1]\).  We choose the usual bounded versions, so
\(0\le V_j(\omega)\le1\) for every \(j\) and every outcome \(\omega\);
replacing arbitrary representatives on their null exceptional sets does not
change their joint law or independence.  Let \(U_j=2V_j-1\), so that \(U_j\)
is uniform on \([-1,1]\) and satisfies \(\lvert U_j(\omega)\rvert\le1\)
pointwise.  For \(0<q<1\), define
\begin{align}
 Y_q&=(1-q)\sum_{j=0}^{\infty}q^jV_j,\label{eq:fr-Y}\\
 X_q&=2Y_q-1=(1-q)\sum_{j=0}^{\infty}q^jU_j.\label{eq:fr-X}
\end{align}
The first variable takes values in \([0,1]\), while the centered variable
takes values in \([-1,1]\).  Write \(F_q(x)=\Probability(Y_q\le x)\).

\begin{theorem}[Geometric law, refinement, and smooth density]
\label{thm:fr-geometric-density}
The series in \cref{eq:fr-Y,eq:fr-X} converge absolutely for every outcome.
The characteristic function of \(X_q\) is
\begin{equation}
 \phi_q(t)=\Expectation\EulerE^{\ImaginaryUnit tX_q}
 =\prod_{j=0}^{\infty}\SincRad\bigl((1-q)q^jt\bigr),
 \qquad \SincRad z=\frac{\sin z}{z},
 \label{eq:fr-sinc-product}
\end{equation}
with the removable value \(\SincRad 0=1\); equivalently,
\(\SincPi x=\SincRad(\pi x)\).  For every \(A>0\) there is a constant
\(C_{A,q}\) such that
\begin{equation}
 |\phi_q(t)|\le C_{A,q}(1+|t|)^{-A}.
 \label{eq:fr-fourier-rapid}
\end{equation}
Consequently \(X_q\) has a compactly supported \(C^\infty\) density \(u_q\),
and \(F_q\) is \(C^\infty\).  The distributional decomposition
\[
 X_q\EqualInLaw (1-q)U_0+qX_q',
\]
where \(X_q'\) is an independent copy, gives
\begin{align}
 u_q(x)
 &=\frac1{2(1-q)}
   \int_{(x-(1-q))/q}^{(x+(1-q))/q}u_q(s)\Differential s,
 \label{eq:fr-density-refinement}\\
 u_q'(x)
 &=\frac{u_q((x+(1-q))/q)-u_q((x-(1-q))/q)}
 {2q(1-q)}.\label{eq:fr-density-derivative}
\end{align}
Equivalently, the CDF of \(Y_q\) satisfies
\begin{align}
 F_q(x)
 &=\int_0^1F_q\left(\frac{x-(1-q)v}{q}\right)\Differential v,
 \label{eq:fr-cdf-refinement}\\
 F_q'(x)
 &=\frac{F_q(x/q)-F_q((x-(1-q))/q)}{1-q}.
 \label{eq:fr-cdf-derivative}
\end{align}
\end{theorem}

\begin{proof}
Since the nonnegative weights \((1-q)q^j\) sum to one and
\(|U_j|\le1\), both random series converge absolutely and have the asserted
supports.  For the \(N\)-th partial sum, independence and
\(\Expectation\EulerE^{\ImaginaryUnit sU_j}=\SincRad s\) give the finite version of
\cref{eq:fr-sinc-product}.  The partial sums converge pointwise and are
bounded in modulus by one, so dominated convergence gives the infinite
product.

For real \(s\), \(|\SincRad s|\le\min(1,|s|^{-1})\).  Fix an integer \(r>A\).
Once \(|t|\ge((1-q)q^{r-1})^{-1}\), the first \(r\) factors therefore give
\[
 |\phi_q(t)|
 \le \prod_{j=0}^{r-1}\frac1{(1-q)q^j|t|}
 =\frac{q^{-r(r-1)/2}}{(1-q)^r}|t|^{-r}.
\]
The remaining factors have modulus at most one.  Enlarging the constant on a
bounded interval proves \cref{eq:fr-fourier-rapid}.  Hence
\(t^m\phi_q(t)\) is integrable for every \(m\).  Fourier inversion and
differentiation under the integral produce a \(C^\infty\) density for the
 law of \(X_q\); its support is already contained in \([-1,1]\).  Since
 \(Y_q=(X_q+1)/2\), its density is \(2u_q(2x-1)\), and hence
 \(F_q'(x)=2u_q(2x-1)\).  Thus \(F_q\) is \(C^\infty\) as claimed.

Removing the first summand in \cref{eq:fr-X} leaves \(q\) times an independent
copy.  Conditioning on either summand gives
\cref{eq:fr-density-refinement}; the fundamental theorem of calculus gives
\cref{eq:fr-density-derivative}.  The same conditioning applied to
\cref{eq:fr-Y} gives \cref{eq:fr-cdf-refinement}.  Changing variables in that
integral yields
\[
 F_q(x)=\frac q{1-q}
 \int_{(x-(1-q))/q}^{x/q}F_q(s)\Differential s,
\]
whose derivative is \cref{eq:fr-cdf-derivative}.
\end{proof}

\section{The exact Fabius and Rvachev member}
\label{sec:fr-fabius-member}

\begin{definition}[Normalized bounded Fabius function]
\label{def:fr-normalized-fabius}
A normalized bounded Fabius function is a continuously differentiable map
\(G:\RealNumbers\to[0,1]\) such that
\[
 G(x)=0\quad(x\le0),\qquad G(x)=1\quad(x\ge1),
 \qquad G(1-x)=1-G(x),
\]
and
\[
 G'(x)=2G(2x)\qquad(0\le x\le\tfrac12).
\]
This exterior-tail normalization removes the additive and scaling ambiguities
from the differential equation.
\end{definition}

\begin{theorem}[The half-base CDF is the Fabius function]
\label{thm:fr-fabius-identification}
There is a unique normalized bounded Fabius function in the sense of
\cref{def:fr-normalized-fabius}; denote it by \(\FabiusBounded\).  Then
\begin{equation}
 F_{1/2}(x)=\FabiusBounded(x)\qquad(x\in\RealNumbers).
 \label{eq:fr-fabius-cdf}
\end{equation}
Moreover \(u_{1/2}\) is Rvachev's compactly supported function:
\begin{equation}
 u_{1/2}(s)=\RvachevUp(s)=
 \begin{cases}
   \FabiusBounded(s+1),&-1\le s\le0,\\
   \FabiusBounded(1-s),&0\le s\le1,\\
   0,&|s|\ge1.
 \end{cases}
 \label{eq:fr-up-fold}
\end{equation}
With the Fourier convention
\(\FourierTwoPi{u}(z)=\int_\RealNumbers u(s)\EulerE^{-2\pi\ImaginaryUnit sz}\Differential s\), for every
\(z\in\RealNumbers\),
\begin{equation}
 \FourierTwoPi{\RvachevUp}(z)=\prod_{j=0}^{\infty}
 \SincPi\left(\frac{z}{2^j}\right).
 \label{eq:fr-rvachev-product}
\end{equation}
Thus the CDF statement and the translated-density statement are compatible;
they are not competing normalizations.
\end{theorem}

\begin{proof}
The law of \(Y_{1/2}\) is contained in \([0,1]\), and
\cref{thm:fr-geometric-density} gives it a smooth density; in particular it
has no atoms.  Replacing every \(V_j\) by \(1-V_j\) shows that \(Y_{1/2}\)
and \(1-Y_{1/2}\) have the same law.  Therefore
\[
 F_{1/2}(1-x)=\Probability(Y_{1/2}\ge x)=1-F_{1/2}(x),
\]
where atomlessness justifies the second equality.  When
\(0\le x\le1/2\),
\cref{eq:fr-cdf-derivative} becomes
\[
 F_{1/2}'(x)=2\bigl(F_{1/2}(2x)-F_{1/2}(2x-1)\bigr)
 =2F_{1/2}(2x),
\]
because \(2x-1\le0\).  Together with the two constant exterior tails and the
reflection law, these are precisely the normalized bounded Fabius axioms.

For completeness, uniqueness follows from a contraction argument.  Any
candidate \(G\) in \cref{def:fr-normalized-fabius} is nondecreasing: its
derivative is nonnegative on \([0,1/2]\), and differentiating reflection gives
\(G'(1-x)=G'(x)\).  Thus \(G\) is the CDF of a probability measure \(\mu\)
on \([0,1]\).  For such a CDF define
\[
 (\mathcal TG)(x)=\int_0^1G(2x-v)\Differential v
 =\int_{2x-1}^{2x}G(s)\Differential s.
\]
This is the CDF of \((V+Z)/2\), where \(V\) is uniform on \([0,1]\),
\(Z\) has CDF \(G\), and they are independent.  On \((0,1/2)\),
\[
 (\mathcal TG)'(x)=2G(2x)=G'(x),
\]
while on \((1/2,1)\), reflection and its differentiated form give
\[
 (\mathcal TG)'(x)=2\bigl(1-G(2x-1)\bigr)
 =2G(2-2x)=G'(1-x)=G'(x).
\]
Both functions vanish at zero and have the same constant exterior tails, so
\(\mathcal TG=G\): every normalized bounded Fabius law is a fixed point.

On probability measures on \([0,1]\), let \(\mathcal T\mu\) denote the law
just constructed.  Given any coupling \((Z_1,Z_2)\) of \(\mu_1,\mu_2\), use
the same independent \(V\) in both images.  Then
\[
 \Expectation\left|\frac{V+Z_1}{2}-\frac{V+Z_2}{2}\right|
 =\frac12\Expectation|Z_1-Z_2|,
\]
and taking the infimum over couplings yields
\(W_1(\mathcal T\mu_1,\mathcal T\mu_2)
 \le\tfrac12W_1(\mu_1,\mu_2)\).  Two fixed points therefore have zero
Wasserstein distance.  Finally, separating the first summand shows
\(Y_{1/2}\EqualInLaw(V_0+Y_{1/2}')/2\), so its law is a fixed point.  It is
the unique one, proving \cref{eq:fr-fabius-cdf}.  This is the
human-readable argument whose conclusion is formalized by
\lean{geometricUniformCDF_one_half_eq_fabiusReal}.

Since \(X_{1/2}=2Y_{1/2}-1\), change of variables gives
\[
 u_{1/2}(s)=\frac12F_{1/2}'\left(\frac{s+1}{2}\right).
\]
For \(-1\le s\le0\), the Fabius derivative equation makes the right side
\(\FabiusBounded(s+1)\).  Reflection gives the branch
\(\FabiusBounded(1-s)\) for \(0\le s\le1\), and
the support gives zero outside.  Finally, for real \(z\), the Fourier-sign
convention gives
\(\FourierTwoPi{\RvachevUp}(z)=\phi_{1/2}(-2\pi z)=\phi_{1/2}(2\pi z)\); the second equality
uses the evenness of the centered law.  Substituting \(q=1/2\) in
\cref{eq:fr-sinc-product} and using
\(\SincPi x=\SincRad(\pi x)\) gives \cref{eq:fr-rvachev-product}.
\end{proof}

\section{Thue--Morse spline approximants}
\label{sec:fr-thue-morse}

Let \(\BinaryDigitSum(k)\) be the number of ones in the binary expansion of \(k\), and put
\(\ThueMorseSign{k}=(-1)^{\BinaryDigitSum(k)}\).  The occurrence of these signs is ordinary
inclusion--exclusion on the dyadic geometric comb.

\begin{theorem}[Exact Thue--Morse partial CDF]
\label{thm:fr-thue-morse-splines}
For \(N\ge1\), put
\[
 Y^{(N)}=\sum_{j=0}^{N-1}2^{-j-1}V_j,
 \qquad F_N(x)=\Probability(Y^{(N)}\le x),
 \qquad r_+=\max(r,0).
\]
Then
\begin{equation}
 F_N(x)=\frac{2^{N(N+1)/2}}{N!}
 \sum_{k=0}^{2^N-1}\ThueMorseSign{k}
 \left(x-\frac{k}{2^N}\right)_+^N.
 \label{eq:fr-thue-morse-spline}
\end{equation}
For the bounded Fabius function \(\FabiusBounded\) and every real \(x\),
\begin{equation}
 F_N(x-2^{-N})\le \FabiusBounded(x)\le F_N(x).
 \label{eq:fr-spline-sandwich}
\end{equation}
In particular the explicit splines converge pointwise to \(\FabiusBounded\), and uniformly
with the quantitative bound
\begin{equation}
 0\le F_N(x)-\FabiusBounded(x)
 \le \FabiusBounded(x+2^{-N})-\FabiusBounded(x)
 \le \omega_{\FabiusBounded}(2^{-N}),
 \label{eq:fr-spline-uniform-bound}
\end{equation}
where
\(\omega_{\FabiusBounded}(\delta)=
\sup_{|a-b|\le\delta}|\FabiusBounded(a)-\FabiusBounded(b)|\).  Consequently
\(\|F_N-\FabiusBounded\|_\infty
\le\omega_{\FabiusBounded}(2^{-N})\to0\).
\end{theorem}

\begin{proof}
For positive \(s\), the Laplace transform of a uniform variable on
\([0,w]\) is \((1-\EulerE^{-sw})/(sw)\).  Therefore the Laplace transform, in the
space variable \(x\), of the CDF of a sum of \(N\) such variables is
\[
 \frac1{s^{N+1}\prod_jw_j}\prod_{j=0}^{N-1}(1-\EulerE^{-sw_j}).
\]
Expanding the last product over subsets and using
\[
 \int_0^\infty \EulerE^{-sx}\frac{(x-a)_+^N}{N!}\Differential x
 =\frac{\EulerE^{-sa}}{s^{N+1}}
\]
gives the usual inclusion--exclusion formula.  Here
\(\prod_jw_j=2^{-N(N+1)/2}\).  The subset sums
\(\sum_j\varepsilon_j2^{-j-1}\) are \(k/2^N\) in bit-reversed order, and bit
reversal preserves the number of selected indices.  The subset sign is thus
\(\ThueMorseSign{k}\), proving \cref{eq:fr-thue-morse-spline}.

The omitted tail of \(Y_{1/2}\) lies pointwise, for every outcome, between
zero and \(\sum_{j\ge N}2^{-j-1}=2^{-N}\).  Hence
\[
 \{Y^{(N)}\le x-2^{-N}\}
 \subseteq\{Y_{1/2}\le x\}
 \subseteq\{Y^{(N)}\le x\}.
\]
Taking probabilities and using \cref{eq:fr-fabius-cdf} proves the sandwich.
Its right inequality gives \(\FabiusBounded(x)\le F_N(x)\).  Applying its left inequality
at \(x+2^{-N}\) gives \(F_N(x)\le\FabiusBounded(x+2^{-N})\), and hence
\cref{eq:fr-spline-uniform-bound}.  The continuous function \(\FabiusBounded\) is constant
off \([0,1]\), so it is uniformly continuous on \(\RealNumbers\) and
\(\omega_{\FabiusBounded}(2^{-N})\to0\).  This proves the claimed uniform convergence.
\end{proof}

The finite sign identity behind the theorem is
\begin{equation}
 \prod_{j=0}^{N-1}(1-z^{2^j})
 =\sum_{k=0}^{2^N-1}\ThueMorseSign{k}z^k.
 \label{eq:fr-thue-morse-product}
\end{equation}
It follows immediately by choosing either \(1\) or \(-z^{2^j}\) from each
factor; uniqueness of binary expansion identifies the chosen subset with
\(k\), and its cardinality with \(\BinaryDigitSum(k)\).

\section{Cumulants as exact inverse observables}
\label{sec:fr-cumulant-recovery}

We use the Bernoulli-number convention \(B_2=1/6\), \(B_4=-1/30\).

\begin{theorem}[All cumulants and monotone base recovery]
\label{thm:fr-cumulant-recovery}
For every \(m\ge1\),
\begin{equation}
 \kappa_{2m}(X_q)
 =\frac{2^{2m}B_{2m}}{2m}
   \frac{(1-q)^{2m}}{1-q^{2m}},
 \label{eq:fr-cumulants}
\end{equation}
and every odd cumulant vanishes:
\begin{equation}
 \kappa_{2m+1}(X_q)=0\qquad(m\ge0).
 \label{eq:fr-odd-cumulants}
\end{equation}
After division by the nonzero uniform cumulant, set
\[
 R_m(q)=\frac{\kappa_{2m}(X_q)}{\kappa_{2m}(U_0)}
 =\frac{(1-q)^{2m}}{1-q^{2m}}.
\]
Then \(R_m\) is a strictly decreasing bijection from \((0,1)\) to \((0,1)\).
Thus, given an observed normalized cumulant \(\rho_m\in(0,1)\), the unique
recovered base is equivalently the unique \(q\in(0,1)\) satisfying
\begin{equation}
 \rho_m(1-q^{2m})=(1-q)^{2m}.
 \label{eq:fr-cumulant-polynomial}
\end{equation}
Its endpoint expansions are
\begin{align}
 q&=\frac{1-\rho_m}{2m}+\BigO((1-\rho_m)^2),
 &&\rho_m\uparrow1,\label{eq:fr-cumulant-qzero}\\
 1-q&=(2m\rho_m)^{1/(2m-1)}
 \left(1+\BigO(\rho_m^{1/(2m-1)})\right),
 &&\rho_m\downarrow0.\label{eq:fr-cumulant-qone}
\end{align}
\end{theorem}

\begin{proof}
The moment generating function of \(U_0\) is \(\sinh z/z\).  Near the origin,
differentiating its logarithm gives
\[
 \frac{\Differential}{\Differential z}\log\frac{\sinh z}{z}
 =\coth z-\frac1z
 =\sum_{m\ge1}\frac{2^{2m}B_{2m}}{(2m)!}z^{2m-1}.
\]
After integration, comparison with
\(\log\Expectation\EulerE^{zU_0}=\sum_r\kappa_r(U_0)z^r/r!\) gives
\(\kappa_{2m}(U_0)=2^{2m}B_{2m}/(2m)\) and zero odd cumulants.

Here is the limit argument that justifies adding the countably many
cumulants.  Put
\[
 X_{q,N}=(1-q)\sum_{j=0}^{N}q^jU_j.
\]
The chosen bounded versions give the deterministic estimate
\(\lvert X_q-X_{q,N}\rvert\le q^{N+1}\), while both variables take values in
\([-1,1]\).  Hence their moment generating functions converge locally
uniformly for complex \(z\).  On a sufficiently small fixed disk about zero,
\[
 \bigl|\Expectation\EulerE^{zX_{q,N}}-1\bigr|\le \EulerE^{|z|}-1<1
\]
uniformly in \(N\), and the same estimate holds for \(X_q\).  The analytic
logarithms normalized to vanish at zero therefore converge locally uniformly;
by the Cauchy integral formula, all their derivatives at zero converge.

For each finite \(N\) and \(r\ge1\), independence and homogeneity of cumulants give
\[
 \kappa_r(X_{q,N})
 =\kappa_r(U_0)(1-q)^r\sum_{j=0}^{N}q^{jr}.
\]
Passing to the limit and summing the geometric series proves
\cref{eq:fr-cumulants}.  Symmetry of every finite partial sum, or directly of
the full series, proves \cref{eq:fr-odd-cumulants}, including
\(\kappa_1(X_q)=\Expectation X_q=0\).

Logarithmic differentiation gives
\[
 \frac{R_m'(q)}{R_m(q)}
 =-\frac{2m}{1-q}+\frac{2mq^{2m-1}}{1-q^{2m}}<0,
\]
because
\(q^{2m-1}(1-q)<1-q^{2m}=(1-q)(1+q+\cdots+q^{2m-1})\).
The endpoint limits are one and zero, so strict monotonicity proves the
bijection.  Substituting the observed value \(\rho_m=R_m(q)\) gives
\cref{eq:fr-cumulant-polynomial} and its unique root in \((0,1)\).
At \(q=0\), \(R_m(q)=1-2mq+\BigO(q^2)\).  At \(q=1-\varepsilon\),
\[
 R_m(q)=\frac{\varepsilon^{2m-1}}{2m}
 (1+\BigO(\varepsilon)).
\]
Ordinary and Puiseux reversion give
\cref{eq:fr-cumulant-qzero,eq:fr-cumulant-qone}.
\end{proof}

The case \(m=1\) is especially simple.

\begin{corollary}[Exact variance inverse]
\label{cor:fr-variance-inverse}
The centered variance and its inverse are
\begin{equation}
 V(q)=\Variance(X_q)=\frac{1-q}{3(1+q)},
 \qquad q=\frac{1-3V}{1+3V}.
 \label{eq:fr-variance-inverse}
\end{equation}
The Rvachev value \(q=1/2\) is characterized by \(V=1/9\).
\end{corollary}

\begin{proof}
Either put \(m=1\) in \cref{eq:fr-cumulants}, or sum
\(\frac13(1-q)^2\sum_{j\ge0}q^{2j}\).  Solving the resulting
linear-fractional equation gives the inverse and the specialization.
\end{proof}

\section{Scale-free recovery and an upgraded conjecture}
\label{sec:fr-scale-free}

\begin{theorem}[Kurtosis inverse and total location--scale identifiability]
\label{thm:fr-total-identifiability}
The excess kurtosis of \(X_q\) is
\begin{equation}
 \gamma_2(q)=\frac{\kappa_4(X_q)}{\kappa_2(X_q)^2}
 =-\frac65\frac{1-q^2}{1+q^2},
 \label{eq:fr-kurtosis}
\end{equation}
and is a strictly increasing bijection from \((0,1)\) to \((-6/5,0)\), with
inverse
\begin{equation}
 q=\sqrt{\frac{1+5\gamma_2/6}{1-5\gamma_2/6}}.
 \label{eq:fr-kurtosis-inverse}
\end{equation}
Consequently the mean, positive variance, and excess kurtosis uniquely
identify every member
\[
 Z=\mu+cX_q,
 \qquad \mu\in\RealNumbers,\quad c>0,\quad0<q<1.
\]
Explicitly, \(\mu=\Expectation Z\), \(q\) is given by
\cref{eq:fr-kurtosis-inverse}, and
\begin{equation}
 c=\sqrt{\Variance(Z)\,\frac{3(1+q)}{1-q}}.
 \label{eq:fr-scale-inverse}
\end{equation}
Thus the total cumulant identifiability statement left conjectural in one
source report is a theorem for the stated three-parameter family.
\end{theorem}

\begin{proof}
Insert \(B_2=1/6\) and \(B_4=-1/30\) into
\cref{eq:fr-cumulants}.  Cancellation gives
\[
 \frac{-2}{15}\frac{(1-q)^4}{1-q^4}
 \left(\frac{3(1-q^2)}{(1-q)^2}\right)^2
 =-\frac65\frac{1-q^2}{1+q^2}.
\]
The last rational function is strictly increasing, with the stated endpoint
limits.  Solving it for \(q^2\) gives
\cref{eq:fr-kurtosis-inverse}.  Translation changes only the mean, positive
scaling multiplies variance by \(c^2\), and excess kurtosis is unchanged by
both operations.  Hence the three displayed reconstruction formulas exist
and are unique.  Formula \cref{eq:fr-scale-inverse} is just
\(\Variance(Z)=c^2V(q)\) with \cref{eq:fr-variance-inverse} substituted.
\end{proof}

For redundancy checks, the next standardized cumulant is
\begin{equation}
 \frac{\kappa_6(X_q)}{\kappa_2(X_q)^3}
 =\frac{48}{7}\frac{(1-q^2)^2}{1+q^2+q^4}.
 \label{eq:fr-standardized-sixth}
\end{equation}
At \(q=1/2\), the exact diagnostic triple begins
\[
 \Variance(X_{1/2})=\frac19,
 \qquad\gamma_2(X_{1/2})=-\frac{18}{25},
 \qquad\frac{\kappa_6(X_{1/2})}{\kappa_2(X_{1/2})^3}=\frac{144}{49}.
\]
Equation \cref{eq:fr-standardized-sixth} follows from
\(\kappa_6(U_0)=16/63\) and the same geometric-series cancellation as above.

\section{Endpoint regimes and conditioning}
\label{sec:fr-endpoints}

\begin{proposition}[Uniform and Gaussian endpoint laws]
\label{prop:fr-endpoint-laws}
As \(q\downarrow0\), for the chosen pointwise bounded versions,
\(X_q\to U_0\) uniformly in the sample point (and hence almost surely).  As
\(q\uparrow1\), \(X_q\to0\) in \(L^2\).  After standardization,
\begin{equation}
 \frac{X_q}{\sqrt{V(q)}}\ \ConvergesInLaw\ \mathcal N(0,1).
 \label{eq:fr-central-limit}
\end{equation}
\end{proposition}

\begin{proof}
The first limit has the deterministic bound
\[
 |X_q-U_0|
 \le q|U_0|+(1-q)\sum_{j\ge1}q^j|U_j|\le2q.
\]
The \(L^2\) assertion follows from
\(\Expectation X_q^2=V(q)\to0\).  For the Gaussian limit, put
\[
 a_{q,j}=\frac{(1-q)q^j}{\sqrt{V(q)}}.
\]
Then
\[
 \sum_{j\ge0}a_{q,j}^2=3,
 \qquad \max_{j\ge0}a_{q,j}=a_{q,0}
 =\sqrt{3(1-q^2)}\longrightarrow0.
\]
The sinc product gives the characteristic function of the standardized law as
\(\prod_{j\ge0}\SincRad(t a_{q,j})\).  For fixed real \(t\) and \(q\) sufficiently
close to one, every argument lies in a fixed neighborhood of zero on which
\[
 \log\SincRad z=-\frac{z^2}{6}+\BigO(z^4),
\]
with a uniform remainder constant.  The resulting logarithmic series is
absolutely convergent because \(\sum_j a_{q,j}^2=3\).  Consequently
\[
 \log\prod_{j\ge0}\SincRad(t a_{q,j})
 =-\frac{t^2}{6}\sum_{j\ge0}a_{q,j}^2
  +\BigO\!\left(t^4\sum_{j\ge0}a_{q,j}^4\right)
 =-\frac{t^2}{2}+o(1),
\]
because
\(\sum_j a_{q,j}^4\le(\max_j a_{q,j}^2)\sum_j a_{q,j}^2\to0\).
Thus the characteristic functions converge pointwise to \(\EulerE^{-t^2/2}\),
and L\'evy's continuity theorem proves \cref{eq:fr-central-limit}.  This direct
Fourier argument also avoids applying a finite-row triangular-array theorem to
the countably infinite series.
\end{proof}

Variance recovery is absolutely well-conditioned on every closed
subinterval of \(0<V<1/3\), since
\[
 \frac{\Differential q}{\Differential V}=-\frac6{(1+3V)^2}.
\]
Its relative condition number diverges as \(q\downarrow0\), but that is the
inevitable normalization by a recovered parameter tending to zero, not an
absolute singularity.  Near \(q=1\), the Puiseux behavior in
\cref{eq:fr-cumulant-qone} shows why high-order cumulant recovery is
progressively flatter than variance recovery.

\section{Parameter scores, local CDF inversion, and the symmetry obstruction}
\label{sec:fr-cdf-inverse}

Put \(w_j(q)=(1-q)q^j\).  Then
\[
 w_j'(q)=q^{j-1}\bigl(j-(j+1)q\bigr),
\]
where the formula at \(j=0\) is interpreted by cancellation as \(-1\).
Write
\[
 H_q(x)=\Probability(X_q\le x)=F_q\left(\frac{x+1}{2}\right)
\]
for the centered CDF.

\begin{proposition}[Fourier score and quantile sensitivity]
\label{prop:fr-fourier-score}
Away from the zeros of the sinc factors,
\begin{equation}
 \partial_q\log\phi_q(t)
 =\sum_{j\ge0}w_j'(q)
 \left(t\cot(w_j(q)t)-\frac1{w_j(q)}\right).
 \label{eq:fr-fourier-score}
\end{equation}
The bracket is interpreted by analytic continuation at \(t=0\), where its
value is zero.
The series converges locally uniformly in \(q\in(0,1)\) and real \(t\) on the
open set where none of those factors vanishes.  More quantitatively, for every
compact \(K\Subset(0,1)\), every \(A>0\), and all integers \(k\ge0\), there is
a constant \(C_{A,k,K}\) such that
\begin{equation}
 \sup_{q\in K}\left|\partial_q^k\phi_q(t)\right|
 \le C_{A,k,K}(1+|t|)^{-A}.
 \label{eq:fr-q-derivative-rapid}
\end{equation}
Consequently the characteristic function, density, centered CDF \(H_q\), and
uncentered CDF \(F_q(x)=H_q(2x-1)\) are \(C^\infty\) in \(q\) on compact
subintervals of \((0,1)\), and parameter differentiation may be passed
through Fourier inversion.  If a centered quantile branch satisfies
\(H_q(x_p(q))=p\) and \(u_q(x_p(q))>0\), then it is locally smooth and
\begin{equation}
 x_p'(q)=-\frac{\partial_q\Probability(X_q\le x)|_{x=x_p(q)}}{u_q(x_p(q))}.
 \label{eq:fr-quantile-sensitivity}
\end{equation}
\end{proposition}

\begin{proof}
For a nonzero sinc factor,
\[
 \frac{\Differential}{\Differential w}\log\SincRad(wt)=t\cot(wt)-\frac1w.
\]
At \(t=0\) the apparent singularity is removable because
\[
 t\cot(wt)-\frac1w=-\frac{wt^2}{3}-\frac{w^3t^4}{45}
 +\BigO(w^5t^6).
\]
The chain rule gives each summand in \cref{eq:fr-fourier-score}.  For large
\(j\), the bracket equals
\(-w_jt^2/3+\BigO(w_j^3t^4)\), while
\(|w_j'|\le C_K(j+1)q_+^{j-1}\) on any compact \(K\subset(0,1)\), where
\(q_+=\max K<1\).  The resulting majorant is a convergent multiple of
\(\sum(j+1)q_+^{2j}\), proving local uniform convergence away from the zeros.

For completeness, we quantify the product differentiation, including at its
zeros.  For every \(r\ge0\), direct differentiation of
\(w_j=(1-q)q^j\) gives constants depending only on \(K,r\) such that
\[
 \sup_{q\in K}|\partial_q^r w_j(q)|
 \le C_{K,r}(j+1)^r q_+^j,
 \qquad
 \sum_{j\ge0}\sup_{q\in K}|\partial_q^r w_j(q)|<\infty.
\]
The integral representation
\(\SincRad y=\frac12\int_{-1}^{1}\EulerE^{\ImaginaryUnit yu}\Differential u\) shows that every real
derivative of \(\SincRad\) is bounded.  These two estimates, together with the
chain rule, give normally convergent differentiated products on compact
\((q,t)\)-sets.

Now fix \(k\) and choose an integer \(R>A+2k\).  On \(K\), each of the first
\(R\) weights is bounded away from zero, so an undifferentiated leading factor
is \(\BigO_K(|t|^{-1})\).  In a term of the \(k\)-fold Leibniz rule, suppose
that \(s\le k\) of those leading factors are differentiated.  Their total
derivative order is at most \(k\), so boundedness of the sinc derivatives
makes the term \(\BigO_{K,k,R}(|t|^{k+s-R})\).  Since \(s\le k\) and
\(R>A+2k\), this is \(\BigO_{K,k,R}(|t|^{-A})\).
Derivatives falling on later factors are summable by the preceding weight
estimates and do not change this worst power.  Summing the normally convergent
Leibniz expansion proves \cref{eq:fr-q-derivative-rapid}; enlarging the
constant handles bounded \(t\).

Given any numbers of spatial and parameter derivatives, choose \(A\) large
enough that the corresponding power of \(t\) is integrable.  Dominated Fourier
inversion then proves smooth dependence of \(u_q\).  Since all centered laws
are supported in the common interval \([-1,1]\),
\(H_q(x)=\int_{-1}^xu_q(s)\Differential s\) on that interval, with constant exterior
tails; hence the same conclusion holds for \(H_q\) and \(F_q\).  Finally, the
implicit-function theorem applies at a quantile where \(u_q(x_p)>0\).
Differentiating its defining identity and solving
\(\partial_qH_q(x_p)+u_q(x_p)x_p'=0\) gives
\cref{eq:fr-quantile-sensitivity}.
\end{proof}

For fixed \(x\), the \(C^2\) inverse sensitivity theorem
\cref{thm:uq-inverse-sensitivities} applies to \(q\mapsto F_q(x)\) at any point
where its \(q\)-derivative is nonzero.  If
\(A=\partial_qF_q(x)|_{q=q_0}\), \(B=\partial_q^2F_q(x)|_{q=q_0}\), and
\(\Delta=p-F_{q_0}(x)\), the first terms are
\[
 q-q_0=\frac\Delta A-\frac{B}{2A^3}\Delta^2+\BigO(\Delta^3).
\]
At \(x=1/2\), however, reflection forces \(F_q(1/2)=1/2\) for every \(q\).
The entire parameter score vanishes there, so the midpoint contains no
information about \(q\).  This is an exact identifiability obstruction, not a
numerical conditioning accident.

\section{The geometric comb as a \(q\)-Pochhammer interpolation basis}
\label{sec:fr-geometric-interpolation}

\begin{proposition}[Newton remainder on geometric nodes]
\label{prop:fr-newton-qpochhammer}
Let \(n\ge1\), let \(0<q<1\), let \(I\subset\RealNumbers\) be an interval containing
\(x,1,q,\ldots,q^{n-1}\), and let \(g:I\to\RealNumbers\) be \(n\)-times continuously
differentiable.  Let \(N_{n-1}g\) be its interpolation polynomial at those
\(n\) distinct geometric nodes.  Then some point \(\xi\) in the convex hull of
\(x,1,q,\ldots,q^{n-1}\) satisfies
\begin{equation}
 g(x)-N_{n-1}g(x)
 =\frac{g^{(n)}(\xi)}{n!}
 \prod_{j=0}^{n-1}(x-q^j)
 =\frac{g^{(n)}(\xi)}{n!}x^n\QPochhammer{x^{-1}}{q}{n},
 \label{eq:fr-newton-qpochhammer}
\end{equation}
where the final expression is understood by polynomial continuation at
\(x=0\).
\end{proposition}

\begin{proof}
If \(x=q^j\) for one of the interpolation nodes, both sides of the first
equality vanish, and any \(\xi\) in the convex hull may be chosen.  Otherwise
the \(n+1\) points are distinct.  Put
\[
 c=\frac{g(x)-N_{n-1}g(x)}{\prod_{j=0}^{n-1}(x-q^j)},
 \qquad
 h(t)=g(t)-N_{n-1}g(t)-c\prod_{j=0}^{n-1}(t-q^j).
\]
The real-valued function \(h\) has those \(n+1\) distinct zeros.  Applying
Rolle's theorem \(n\) times gives a point \(\xi\) in their convex hull with
\(h^{(n)}(\xi)=0\).  Since \(N_{n-1}g\) has degree at most \(n-1\) and the
node polynomial is monic of degree \(n\), this identity reads
\(g^{(n)}(\xi)-n!c=0\), proving the first equality.  For \(x\ne0\), the second
equality is the factor-by-factor identity
\[
 x^n\QPochhammer{x^{-1}}{q}{n}
 =x^n\prod_{j=0}^{n-1}(1-x^{-1}q^j)
 =\prod_{j=0}^{n-1}(x-q^j).
\]
The final product is a polynomial in \(x\), so this identity defines the
polynomial continuation of the middle expression at \(x=0\).
\end{proof}

This is the precise interpolation link behind the informal phrase
\emph{geometric comb}.  The same nodes \(q^j\) govern the random-series
weights, the cumulant sums, Newton divided differences, and the finite
\(q\)-Pochhammer basis.  Parameter recovery from Fabius--Rvachev observables is
therefore connected to inverse \(q\)-products structurally, rather than by a
mere resemblance of notation.
